\subsubsection{Integral derivatives away from the discriminant}
\label{subsec:baseline-jets}
\hypertarget{source-baseline-jets}{}

Beyond the primes where the jets give a positive gain, the product of local estimates still requires integrality in every order of differentiation. This applies at the level-free places. We keep the
scale of the curve ideal used in the period. Thus the period
lattice \(\Gamma\) has discriminant
\(g_2(\Gamma)^3-27g_3(\Gamma)^2=1\). Use the basis
\(\lambda a,\lambda(b+\sqrt{-d})/2\) of
Section~\ref{subsec:theta-normalization}, with \(a\) odd.
Write \(E=\C/\Gamma\), \(h=\lambda a/2\) modulo \(\Gamma\),
and \(\theta\) for its even degree-one section normalized by
\(\theta(0)=1\). Its sole zero is \(h\).
Fix a coefficient embedding into \(\mathbb C_p\) and extend it
to a field isomorphism on \(\C\) as in
Section~\ref{subsec:word-measure}. Elliptic functions will
be transported as rational functions on the Weierstrass curve.
All formal identities use the transport on coefficients.

\begin{proposition}[Integral jets at level-free primes]\label{prop:level-free-theta-jets}
Use a discriminant-one curve with period lattice \(\Gamma\),
odd \(a\) in the basis of Section~\ref{subsec:theta-normalization},
and the degree-one even section normalized by \(\theta(0)=1\)
with zero at \(\lambda a/2\bmod\Gamma\).
At \(p\ne d\), on the admissible lattice pairs with the
stated level exclusions and prime-to-\(p\) denominators,
let \(a_{\alpha,w}\) be the tested ordinary derivative of multi-order \(\alpha\) at a coefficient place \(w\). Then
\[
 v_w(a_{\alpha,w})\ge0\qquad\text{for every permitted }\alpha,w.
\]
 The bound is uniform in the order of
the derivative and in the coefficient field. It includes
the primes \(2\) and \(3\) when they split in the chosen CM field.
\end{proposition}

\begin{proof}
\proofstep{Choose a good model with the given invariant differential.}
The CM \(j\)-invariant is an algebraic integer. Consequently
the curve has potentially good reduction. To keep track of
the differential, take a good model after extension and choose
its analytic homothety \(c\) from the fixed lattice. Its
discriminant is \(c^{-12}\), so \(c\) has norm one. The change
of Weierstrass variables with zero translations and unit scaling
returns an integral good model with the original \(g_2,g_3\)
invariants. Over \(\mathbb C_p\) the normalized differential
therefore identifies the formal logarithmic variable with the
transported analytic variable \(z\). Write \(T\) for a good
formal parameter and \(z=l(T)\) for the formal logarithm,
with unit tangent. In particular

\[
  D_z=(l'(T))^{-1}\frac{d}{dT},\qquad D_z\big(R[\![T]\!]\big)\subset R[\![T]\!]
\]
 preserves integral formal series
because the invariant differential has integral unit coefficient
in this parameter. For existence over a finite local extension,
the invariants \(g_2,g_3\) are algebraic by the discriminant
identity and \(j=1728g_2^3\). The analytic scale fixes the
ratios of the invariant differentials. \cite{cm-reciprocity,serre-tate}

\proofstep{Check that the odd characteristic remains visible in the special fiber at two.}

For odd \(p\), \(h\) has nonzero reduction by injectivity on
prime-to-\(p\) torsion. At two we use its particular
characteristic. The ring \(O=\Z[(1+\sqrt{-d})/2]\) splits at two,
since \(d\equiv7\pmod8\). In \(O/2O\) the integers
\(\omega=(1+\sqrt{-d})/2\) and \(\bar\omega\) give the
two complementary idempotents. The two idempotents split the point as
\[
       h=\omega h+\bar\omega h,\qquad
       \omega h\ne0,\quad\bar\omega h\ne0\quad\text{in }E[2].
\]
Indeed the second lattice coordinate of the corresponding lift
is \(a/2\) or \(-a/2\), which excludes membership in \(\Gamma\).

Each idempotent cuts out a height-one summand of
\(E[2^\infty]\), as one sees on the generic Tate module.
The action extends over a good model. The special fiber has
dimension one, so one summand has dimension zero and is
\'{e}tale, and the other has dimension one and is connected.
Reduction is injective on points of the \'etale summand.
Thus the reduction of \(h\) has order two at this place as
well.  It holds for either embedding of \(O\).

\proofstep{Use a split dilation with a unit tangent and norm divisible by the residue prime.}

Suppose that \(p\) splits in \(O\), including \(p=2\) and
\(p=3\) in the selected host (\(d\equiv2\pmod3\)).
Order the components of \(O\otimes\Z_p\) with the differential
embedding first.
At a split prime the two idempotents are
\[
  O\otimes\mathbb Z_p\simeq\mathbb Z_p\times\mathbb Z_p,\qquad e_1=(1,0),\quad e_2=(0,1).
\]
 Choose \(\alpha\in O\) congruent to \((1,0)\)
modulo \(p\), nonzero and of norm \(m>1\). Its action
\([\alpha](T)\) on the good formal curve is integral and
has unit linear coefficient \(A\), while \(p\mid m\).
The quadratic sign of the theta vector at a lattice period
\(v\) is \((-1)^{E_0(\omega v,v)}\), with \(E_0\) its alternating
form, as in Section~\ref{subsec:theta-normalization}.

Replacing
\(v\) by \(\alpha v\) multiplies that integer by \(m\).
The Hermitian automorphy exponents also multiply by \(m\).
The quotient
\[
             R(z)=\theta(\alpha z)/\theta(z)^m
\]
is therefore elliptic. Its divisor is supported on \(h\) and
the preimages of \(h\). None reduces to the identity since
\(\bar h\ne0\) and \([\alpha]\) preserves reduction. Both
\(R\) and its inverse have unit integral expansions at the
identity, with constant term one. For detail, its divisor is
defined over a finite extension of the good reduction field,
so the rational function descends there up to a scalar.
A degree-zero principal divisor of rational points has a
one-dimensional space of defining rational functions, by base
change of sections on the proper curve.

If two such functions \(f,g\) have the prescribed divisor and are normalized at zero, then
\[
  {\rm div}(f/g)=0,\qquad (f/g)(0)=1\quad\Longrightarrow\quad f=g.
\]
 The normalization at
zero fixes the scalar. On a smooth model the only possible
component of the divisor through the origin is vertical.
After dividing by the corresponding power of a uniformizer
one has a unit in the local ring, since that ring is regular
and all divisor orders there vanish. Evaluation at the origin
shows that the original vertical order was zero.

Put \(F(T)=\theta(l(T))\) using the fixed good-model coefficient
embedding into \(\mathbb C_p\) and the normalized \(\theta(0)=1\). The identity
\[
             F([\alpha](T))=F(T)^m R(l(T))
\]
gives integrality of \(F\) from its constant term one. At degree
\(j>0\) the unknown coefficient occurs with multiplier
\(A^j-m\), which is a unit, and all remaining terms are
integral by induction. The same conclusion holds for the
inverse series. This argument uses the split CM action on the
ordinary curve, including in residue characteristic three
where changing to a short integral equation could lose a
factor of three. At three the CM norm identity of \eqref{eq:cm-isogeny} gives the same integral correction with
\(\alpha=\sqrt{-d}\). There the correction

\[
  (\alpha^2-1)b_2/12=-((d+1)/3)b_2/4
\]
 is integral for the
general good equation. \label{cor:baseline-jets-1}

\proofstep{Use the logarithmic derivative of the theta section for primes at least five.}

For \(p\ge5\) with \(p\ne d\), the short equation at this
scale is good. Indeed \(j\) integral gives \(g_2\) integral
and the discriminant identity then gives \(g_3\) integral.
Apply \eqref{eq:cm-isogeny} with \(\alpha=\sqrt{-d}\).
Its kernel is killed by \(d\), hence every nonzero kernel point
has integral affine coordinates. The quadratic-coefficient
identity following \eqref{eq:cm-isogeny} gives \(s_2\) integral
because \(\alpha^2-\alpha\bar\alpha=-2d\) is a unit. The
second logarithmic derivative of the even theta function is
\[
           (\log\theta)''(z)=-\wp(z+h)-s_2 .
\]
The shift by \(h\) and the scalar gauge in the translated sigma
function contribute only a linear term to the logarithm in
addition to the displayed quadratic correction. The translated
Weierstrass coordinate is regular integral on the identity
disc by good reduction and \(\bar h\ne0\). All of its invariant
derivatives at the origin are integral.

Since \(\theta(0)=1\)
and \(\theta'(0)=0\), every positive derivative of
\(\log\theta\) at zero is integral. Applying Leibniz's rule
successively to \(\theta'=\theta(\log\theta)'\) proves
integrality of every ordinary derivative of \(\theta\).
This also treats inert primes where the argument with the
split dilation is unavailable.

\proofstep{Remove the zero of a translated section before taking a prime-to-residue-characteristic root.}

Let \(t\) have order dividing \(M\) prime to \(p\). With the
normalized translation \(U_t\) above, put
\[
       Q_t(z)=\left(\frac{U_t\theta(z)}{\theta(z)}\right)^M .
\]
The ratio has a period character of order dividing \(M\) by
the Heisenberg commutator. Thus \(Q_t\) is rational on the
curve and its divisor is \(M([h-t]-[h])\).
Let \(|\cdot|_{\rm vert}\) be the norm at the
generic point of the good special fibre. The telescope
of theta ratios gives
\[
          \prod_{j=0}^{M-1} Q_t(z+jt)\ \in\ \mu_\infty:
\]
after cycling the translates the only factor is a
Heisenberg phase and the \(Mt\)-period sign. Each
translation preserves the special fibre and its generic
valuation. Hence \(|Q_t|_{\rm vert}^M=1\) and
\(|Q_t|_{\rm vert}=1\), including any scalar of the
normalized theta ratio by multiplicativity.

At \(p=2\) the reduction of \(h\) has order two and \(t\)
has odd order. At the odd primes the points generated by
\(h,t\) are prime-to-\(p\) torsion and reduction is injective
on them. Thus the divisor on the identity disc has either
no support or a zero of order \(M\) at the identity, the
latter precisely when \(t=h\) on the curve. The order on the identity disc can thus be written
\[
  \operatorname{ord}_{T=0}(Q_t)=eM,\qquad e=\begin{cases}1&t=h,\\0&t\ne h.\end{cases}
\]
 The expansion of \(Q_t/T^{eM}\)
is an integral unit with unit constant term. To see this, take
a defining function over a finite extension up to a scalar,
remove its vertical order, then remove the power of a regular
parameter defining the identity section.

The remaining
function and inverse are units at the special point by the
local divisor criterion above. The original generic norm
one makes the residual scalar a unit. The formal root
\((U_t\theta/\theta)(l(T))/T^e\) thus has unit constant and
integral coefficients. After dividing by its constant term,
write the normalized root as \(1+\sum_{j\ge1}b_j T^j\).
Then
\[
 [T^k]\Bigl(1+\sum_{j\ge1}b_j T^j\Bigr)^M
      =M b_k+P_k(b_1,\ldots,b_{k-1})
             \quad(P_k\in\mathbb Z[b_1,\ldots,b_{k-1}]).
\]
Since \(M\) is a unit, induction gives integral \(b_k\)
at every positive degree. The ratio itself is an integral series in \(T\).
Consequently its invariant derivatives are integral. The
product rule and the preceding bounds for \(\theta\) prove
the unit bound for all jets of \(U_t\theta\). In particular
the calculation covers the translate whose constant term
vanishes without dividing by that value.

\proofstep{Transfer the unit bound through both integral lattice comparisons.}

Consider a kernel term between the actual source and target
lattices. Its admissible pair agrees locally with these
lattices at \(p\). Choose positive integers \(s_1,s_2\)
prime to \(p\) clearing their two lattice comparisons in both
directions. For every test matrix \(y\) mapping the source
into the target, \(s_1s_2 y\) maps the admissible source into
the admissible target.
Neither clearing factor affects the valuation of a coordinate of a test direction, since
\[
  v_p(s_1s_2)=0,\qquad v_p(s_1s_2v)=v_p(v).
\]
 Multiplication by any period of
\(\Gamma\) sends this scaled embedded direction into the
period lattice of the Hom space. Likewise a translate
prime-to-\(p\) on the original pair has a prime-to-\(p\)
denominator on the admissible pair, by the same inclusions.

Apply the standardizing isometry \(g\) of
Section~\ref{subsec:theta-normalization} to the four copies.
It carries their lattice into the standard one after multiplication
by an integer prime to \(p\), by the relative index.
For each coordinate \(v\) of a scaled test direction the same
integer \(b_0\), prime to \(p\), gives
\[
 b_0v\Gamma\subseteq\Gamma
 \quad\Longrightarrow\quad b_0v\in\operatorname{End}(\Gamma)=O
 \quad\Longrightarrow\quad v\in O\otimes\Z_{(p)}.
\]
Here the endomorphisms are complex multipliers of the period lattice.
This proves integrality of the argument matrix under the chosen
embedding. Undoing the two lattice comparisons only multiplies
its columns by \(1/(s_1s_2)\). Further integral linear changes
of the two local bases, including the dual basis on the source,
preserve this conclusion by the chain rule.

In the four copies the normalized section is a finite sum of
products of the degree-one translates just treated, with
integral scalars. Integral ordinary derivatives are stable under
such sums, products and the indicated linear substitutions.
For instance the product step uses the integral coefficients in
\[
  \partial^\alpha(FG)(0)=\sum_{\beta\le\alpha}\binom{\alpha}{\beta}\partial^\beta F(0)\,\partial^{\alpha-\beta}G(0).
\]
The derivative identity in Lemma~\ref{lem:cm-disjoint-tensor}, applied
to the normalized section with \(m=4\), gives the unit bound for each
ordinary derivative in one copy. A prescribed
integral linear combination of translates satisfies that
bound as well. Use the stated admissible lattices and level exclusions and the prime-to-\(p\) denominator throughout. The estimate has constant one independently
of derivative order, the number of terms, the word and the
degree of the coefficient field. \label{cor:baseline-jets-2}
\end{proof}

